\section*{Chapter \ref{ch:g12:proton-nmr} --- Proton NMR}

\begin{solution}{exo:g12:proton-nmr:1}
Methane: 1. Ethane: 1 (the two \ce{CH3} are alike). Propane: 2 (two
\ce{CH3} alike, one \ce{CH2}). Methanol: 2 (\ce{CH3} and \ce{OH}).
Methoxymethane: 1.
\end{solution}

\begin{solution}{exo:g12:proton-nmr:2}
2 neighbours: a triplet, 1:2:1. 3 neighbours: a quartet, 1:3:3:1.
0 neighbours: a singlet.
\end{solution}

\begin{solution}{exo:g12:proton-nmr:3}
Total \qty{40}{mm} for 8 \omterm{def:g9:inside-the-atom:proton}{protons}: \qty{5}{mm} per \omterm{def:g9:inside-the-atom:proton}{proton}. The signals
represent 3, 2 and 3 \omterm{def:g9:inside-the-atom:proton}{protons}.
\end{solution}

\begin{solution}{exo:g12:proton-nmr:4}
\omterm{def:g11:functional-groups:carbonyl}{Aldehyde}: 9.7 to \qty{10.0}{ppm}. A \ce{CH3} of a chain: 0.7 to
\qty{1.3}{ppm}.
\end{solution}

\begin{solution}{exo:g12:proton-nmr:5}
Its twelve \omterm{def:g12:proton-nmr:equivalent}{equivalent protons} give a single, sharp, intense line, easy
to place at zero. The deuterated \omterm{def:g6:solutions-and-solubility:solute}{solvent} gives no \omterm{def:g9:inside-the-atom:proton}{proton} signal of its
own, which would otherwise swamp those of the few milligrams of sample.
\end{solution}

\begin{solution}{exo:g12:proton-nmr:6}
Two signals. \ce{CH2} next to chlorine: 2 \omterm{def:g9:inside-the-atom:proton}{protons}, a quartet (3
neighbours), in the range 2.5--\qty{4.0}{ppm}. \ce{CH3}: 3 \omterm{def:g9:inside-the-atom:proton}{protons}, a
triplet (2 neighbours), near the usual \ce{CH3} range, a little higher
because the chlorine is two \omterm{def:g7:atoms-and-molecules:atom}{atoms} away.
\end{solution}

\begin{solution}{exo:g12:proton-nmr:7}
Three signals: the two \ce{CH3}, 6 \omterm{def:g12:proton-nmr:equivalent}{equivalent protons}, a doublet (1
neighbour, the \ce{CH}); the \ce{CH}, 1 \omterm{def:g9:inside-the-atom:proton}{proton}, split by 6 neighbours
into 7 lines, in the \ce{H-C-O} range (3.3--\qty{4.5}{ppm}); the
\ce{OH}, 1 \omterm{def:g9:inside-the-atom:proton}{proton}, a singlet.
\end{solution}

\begin{solution}{exo:g12:proton-nmr:8}
(a) propanone; (b) ethyl ethanoate; (c) propanoic acid.
\end{solution}

\begin{solution}{exo:g12:proton-nmr:9}
The \ce{CH2} is bonded to the oxygen \omterm{def:g7:atoms-and-molecules:atom}{atom} of the \omterm{def:g11:functional-groups:carboxylic-acid}{ester}, which is
electronegative and deshields its \omterm{def:g9:inside-the-atom:proton}{protons}: \ce{H-C-O} range,
3.3--\qty{4.5}{ppm}. The \ce{CH3} is one carbon further away and stays in
the chain range.
\end{solution}

\begin{solution}{exo:g12:proton-nmr:10}
Propanone: a \ce{C=O} (\omterm{def:g11:functional-groups:carbonyl}{ketone}, \qty{1715}{cm^{-1}}) and six \omterm{def:g12:proton-nmr:equivalent}{equivalent
protons}, one singlet. Propanal would give three signals, one of them,
the \omterm{def:g11:functional-groups:carbonyl}{aldehyde} \omterm{def:g9:inside-the-atom:proton}{proton} of \ce{CH3-CH2-CHO}, near 9.7--\qty{10.0}{ppm}.
\end{solution}

\begin{solution}{exo:g12:proton-nmr:11}
The \ce{OH} \omterm{def:g9:inside-the-atom:proton}{proton} is exchanged quickly between \omterm{def:g7:atoms-and-molecules:molecule}{molecules}: it gives a
singlet and does not split its neighbours. The \ce{CH2} is therefore
split only by the 3 \omterm{def:g9:inside-the-atom:proton}{protons} of the \ce{CH3}: a quartet.
\end{solution}

\begin{solution}{exo:g12:proton-nmr:12}
Butanoic acid \ce{CH3-CH2-CH2-COOH}: 4 signals, a triplet (\ce{CH3}), a
6-line signal (central \ce{CH2}, 5 neighbours), a triplet
(\ce{CH2-C=O}), a singlet near 11--\qty{12}{ppm}. Ethyl ethanoate: a
singlet, a quartet, a triplet. Methyl propanoate \ce{CH3-CH2-CO-O-CH3}:
a triplet, a quartet (\ce{CH2-C=O}, 2.0--\qty{2.4}{ppm}), a singlet
(\ce{O-CH3}). Propyl methanoate \ce{H-CO-O-CH2-CH2-CH3}: 4 signals, a
singlet far to the left (the H on the \ce{C=O} carbon), a triplet
(\ce{O-CH2}), a 6-line signal, a triplet (\ce{CH3}).
\end{solution}

\begin{solution}{exo:g12:proton-nmr:13}
The \ce{OH} signal: the \ce{O-H} \omterm{def:g9:inside-the-atom:proton}{protons} exchange with the deuterium of
the heavy water, and \ce{O-D} gives no \omterm{def:g9:inside-the-atom:proton}{proton} signal. A signal that
vanishes on shaking with \ce{D2O} belongs to a \omterm{def:g9:inside-the-atom:proton}{proton} on O or N.
\end{solution}

\begin{solution}{exo:g12:proton-nmr:14}
Ethanol adds a quartet near \qty{3.69}{ppm}, a triplet near
\qty{1.23}{ppm} and an \ce{OH} singlet. The \omterm{def:g11:functional-groups:carboxylic-acid}{ester} singlet: \qty{30}{mm}
for 3 \omterm{def:g9:inside-the-atom:proton}{protons}, \qty{10}{mm} per \omterm{def:g9:inside-the-atom:proton}{proton} of \omterm{def:g11:functional-groups:carboxylic-acid}{ester}. The ethanol triplet:
\qty{3}{mm} for 3 \omterm{def:g9:inside-the-atom:proton}{protons}, \qty{1}{mm} per \omterm{def:g9:inside-the-atom:proton}{proton} of ethanol. Since each
\omterm{def:g7:atoms-and-molecules:molecule}{molecule} contributes one \ce{CH3} to these signals, the amounts are in
the ratio $1/10$: about one \omterm{def:g7:atoms-and-molecules:molecule}{molecule} of ethanol for ten of \omterm{def:g11:functional-groups:carboxylic-acid}{ester}.
\end{solution}

\begin{solution}{exo:g12:proton-nmr:15}
Four signals: the two \ce{CH3} (6 \omterm{def:g9:inside-the-atom:proton}{protons}) a doublet; the \ce{CH} (1
\omterm{def:g9:inside-the-atom:proton}{proton}), split by 6 + 2 = 8 neighbours into 9 lines; the \ce{CH2} (2
\omterm{def:g9:inside-the-atom:proton}{protons}) next to \ce{O}, a doublet, in the \ce{H-C-O} range; the
\ce{OH} (1 \omterm{def:g9:inside-the-atom:proton}{proton}) a singlet. The \ce{CH} signal has the most lines.
\end{solution}

\begin{solution}{pb:g12:proton-nmr:1}
\textbf{1.} An \omterm{def:g11:organic-skeletons:alkane}{alkane} with four carbons is \ce{C4H10}; one \ce{C=O}
removes two hydrogens and the two oxygens add none: \ce{C4H8O2}, the
formula of acids and \omterm{def:g11:functional-groups:carboxylic-acid}{esters} with four carbons.

\textbf{2.} Butanoic acid \ce{CH3-CH2-CH2-COOH}; ethyl ethanoate
\ce{CH3-CO-O-CH2-CH3}; methyl propanoate \ce{CH3-CH2-CO-O-CH3}; propyl
methanoate \ce{H-CO-O-CH2-CH2-CH3}.

\textbf{3.} The three \omterm{def:g11:functional-groups:carboxylic-acid}{esters} share the \omterm{def:g11:functional-groups:carboxylic-acid}{ester} group; all four have a
\ce{C=O}.

\textbf{4.} A \ce{C=O}, in the \omterm{def:g11:functional-groups:carboxylic-acid}{ester} range (near \qty{1735}{cm^{-1}}).

\textbf{5.} The very broad \ce{O-H} band from 2500 to
\qty{3300}{cm^{-1}}: absent, so not butanoic acid.

\textbf{6.} No: the three \omterm{def:g11:functional-groups:carboxylic-acid}{esters} have the same groups.

\textbf{7.} Three.

\textbf{8.} An ethyl group \ce{CH3-CH2-}: the \ce{CH2} split by 3
\omterm{def:g9:inside-the-atom:proton}{protons} (quartet), the \ce{CH3} by 2 (triplet).

\textbf{9.} Its \omterm{def:g9:inside-the-atom:proton}{protons} have no \omterm{def:g9:inside-the-atom:proton}{protons} on the neighbouring \omterm{def:g7:atoms-and-molecules:atom}{atoms}: a
\ce{CH3} bonded to the \ce{C=O} carbon or to the \omterm{def:g11:functional-groups:carboxylic-acid}{ester} oxygen.

\textbf{10.} The \ce{CH2} is bonded to oxygen (\ce{H-C-O} range).

\textbf{11.} A triplet (\ce{CH3}), a quartet near 2.0--\qty{2.4}{ppm}
(\ce{CH2} next to \ce{C=O}) and a singlet in the \ce{H-C-O} range
(\ce{O-CH3}). The quartet would be far below \qty{4.12}{ppm} and the
singlet far above \qty{2.04}{ppm}: not the unknown.

\textbf{12.} Four signals (a singlet far left, a triplet, a 6-line
signal, a triplet): not the unknown, which has three.

\textbf{13.} Ethyl ethanoate, \ce{CH3-CO-O-CH2-CH3}: singlet
\qty{2.04}{ppm} = \ce{CH3-C=O}; quartet \qty{4.12}{ppm} = \ce{O-CH2};
triplet \qty{1.26}{ppm} = the end \ce{CH3}.

\textbf{14.} Singlet 3, quartet 2, triplet 3.

\textbf{15.} $72 / 8 = \qty{9}{mm}$ per \omterm{def:g9:inside-the-atom:proton}{proton}.

\textbf{16.} Singlet $3 \times 9 = \qty{27}{mm}$; triplet
\qty{27}{mm}.

\textbf{17.} $18 + 27 + 27 = \qty{72}{mm}$.

\textbf{18.} Yes: small \omterm{def:g11:functional-groups:carboxylic-acid}{esters} are known for their fruity smells.

\textbf{19.} The quartet's step is \qty{18}{mm} out of \qty{72}{mm}.
\end{solution}
